\documentclass[11pt,a4paper]{article}
\usepackage[utf8]{inputenc}
\usepackage[T1]{fontenc}
\usepackage{lmodern}
\usepackage[margin=2.25cm]{geometry}
\usepackage{amsmath,amssymb,mathtools,amsthm,stmaryrd,booktabs,longtable,array,enumitem,xcolor,hyperref,listings}
\hypersetup{colorlinks=true,linkcolor=blue,urlcolor=blue,citecolor=blue}
\emergencystretch=2em
\setlist{leftmargin=1.6em,itemsep=2pt,topsep=4pt}
\newcommand{\code}[1]{\path{#1}}
\newcommand{\kw}[1]{\textsf{#1}}
\newcommand{\Comp}{\mathcal{C}}
\newcommand{\Prog}{\mathcal{P}}
\newcommand{\Eff}{\mathcal{E}}
\newcommand{\Trace}{\mathcal{T}}
\newcommand{\State}{\mathcal{S}}
\newcommand{\Pre}{\mathsf{Pre}}
\newcommand{\Post}{\mathsf{Post}}

\newcommand{\trace}{\mathcal{T}}
\theoremstyle{definition}
\newtheorem{definition}{Définition}[section]
\newtheorem{invariant}[definition]{Invariant}
\theoremstyle{plain}
\newtheorem{proposition}[definition]{Proposition}
\newtheorem{theorem}[definition]{Théorème}
\theoremstyle{remark}
\newtheorem{remark}[definition]{Remarque}
\title{Théorie de la programmation compositionnelle\
\large Fondements, sémantique, preuves, synthèse, résilience et protocole d'évaluation}
\author{Dossier théorique préparatoire à FulleGenBrain\_2}
\date{25 octobre 2026}
\begin{document}
\maketitle
\begin{abstract}
Ce document développe une théorie opératoire de la programmation compositionnelle pour servir de fondation à FulleGenBrain\_2. Le principe est de construire des comportements et des programmes à partir de composants typés, versionnés, observables et vérifiables, au moyen d'une syntaxe compositionnelle déclarative dont la sémantique, les effets, les coûts, les dépendances et les conditions d'échec sont explicités. Le périmètre va des fonctions pures aux compositions de processus avec état, effets, observation, recherche, récupération, vérification, apprentissage et transformation contrôlée de programmes. La théorie distingue les lois prouvées sous hypothèses formelles, les invariants de conception, les mécanismes observés dans les sources AAM O94.6, et les propriétés qui exigent encore des preuves machine ou des expériences plus larges. Elle est conçue pour être extensible : une grammaire ou un type de composant insuffisant est une obligation de formalisation, pas un motif pour abandonner une tâche.
\end{abstract}
\newpage\setcounter{tocdepth}{1}\tableofcontents\newpage

\section{Mandat, terminologie et statut des affirmations}
La programmation compositionnelle est ici une méthode générale permettant de construire une capacité complexe à partir d'unités plus petites dont les contrats et les relations sont explicités. Elle ne se réduit ni à chaîner des fonctions, ni à une syntaxe DSL, ni au fait de ranger des modules dans un catalogue. Une composition complète doit préciser au moins : (i) la structure du programme ; (ii) le typage des flux de données ; (iii) les préconditions et postconditions ; (iv) les effets et ressources ; (v) la sémantique d'exécution ; (vi) la gestion des défaillances ; (vii) la provenance ; (viii) les critères de validation ; (ix) les règles de transformation et de révision.

Les statuts utilisés dans le document sont : \kw{SOURCE-OBSERVED} (élément retrouvé dans les artefacts examinés), \kw{DEFINITION} (construction mathématique donnée ici), \kw{PROVED-UNDER-ASSUMPTIONS} (preuve déductive avec hypothèses identifiées), \kw{DESIGN-INVARIANT} (règle de conception à faire respecter par l'implémentation), \kw{EMPIRICAL-TO-TEST} (effet attendu mais non démontré), \kw{NOT-CHECKED-BY-PROVER} (preuve humaine structurée, pas encore vérifiée par assistant de preuve).

\section{Examen des artefacts AAM O94.6 et portée des constats}
\subsection{Éléments effectivement observés}
L'audit livré dans le paquet AAM O94.6 décrit une séparation entre le code du modèle, les preuves, les reçus, les journaux et les oracles de tests. Le démonstrateur compositionnel joint comprend un catalogue de briques et un DSL JSON avec des nœuds \code{invoke}, \code{sequence}, \code{condition} et \code{map}; les expressions de données prennent notamment la forme de variables, littéraux, listes, objets et opérateurs purs. Le compilateur vérifie les identifiants de composants, la forme des nœuds, la profondeur et le nombre de nœuds, produit un fichier source et un bytecode, puis écrit un manifeste avec empreintes. Le runtime résout les expressions et n'autorise que les identifiants de composants de son catalogue. La composition d'exemple orchestre validation de requête, inventaire de projet, compilation, tests et attachement d'intégrité.

Le paquet de démonstration examiné dans ce travail contient un mini-projet d'éditeur et un ensemble local de cinq tests. Dans l'environnement présent, l'appel initial des tests depuis la mauvaise racine a produit une erreur d'import de module; après lancement depuis le répertoire du projet avec son chemin d'import local, les cinq tests ont réussi. Ce fait valide uniquement les cinq tests de ce mini-projet et non la totalité d'AAM, ni la théorie ici élaborée. Le reçu AAM de la matrice P8 rapporte séparément 92 fichiers de test avec 449 réussites cumulées, zéro échec et zéro timeout dans son environnement d'audit; ce résultat est attribué à ce reçu et à ses limites, il n'est pas présenté comme une exécution effectuée dans le présent travail.

\subsection{Ce que ces constats autorisent à conclure}
Ils fournissent un point d'appui pratique pour la composition déclarative, le registre de composants, la compilation, les manifests, le traçage et les portes de qualité. Ils ne suffisent pas à démontrer l'exhaustivité de la théorie, la correction générale du compilateur, la sûreté de tout composant, la complétude de la synthèse ou la généralisation à un domaine arbitraire. Le présent document formalise donc les obligations mathématiques et d'ingénierie qui sont nécessaires pour élargir ce noyau.

\section{Objet mathématique central}
\subsection{Composant typé}
\begin{definition}[Contrat d'un composant]
Un composant est le tuple
\[
c=(id,ver,A,B,P,Q,E,R,\delta,\pi,V),
\]
où $id$ est l'identifiant stable, $ver$ la version, $A$ et $B$ les types d'entrée et de sortie, $P$ la précondition, $Q$ la postcondition, $E$ l'ensemble des effets possibles, $R$ le profil de ressources, $\delta$ la sémantique d'exécution, $\pi$ le schéma de provenance et $V$ le vérificateur ou protocole d'évaluation associé. Les types peuvent être algébriques, dépendants d'un contexte, structurés ou enrichis de contraintes d'unités et d'effets.
\end{definition}

Un contrat doit distinguer l'interface syntaxique de la garantie sémantique. Déclarer qu'un composant transforme $A$ en $B$ ne prouve pas qu'il le fait toujours. Un niveau de garantie est enregistré dans les métadonnées : spécifié, compilé, testé localement, vérifié formellement sous hypothèses, ou observé en production. Ces niveaux ne se substituent pas les uns aux autres.

\subsection{Types de résultat}
Une exécution renvoie non seulement une valeur, mais un résultat typé :
\[
Y = \mathsf{Success}(v,w) \mid \mathsf{Partial}(v,w) \mid \mathsf{Unresolved}(U,w)
\mid \mathsf{Rejected}(r,w) \mid \mathsf{Failed}(e,w) \mid \mathsf{Cancelled}(w),
\]
où $w$ est une trace de provenance, $U$ un ensemble structuré d'alternatives, $r$ une raison de refus et $e$ une description d'erreur. Une chaîne ne doit pas convertir silencieusement une sortie partielle ou non résolue en succès. Les résultats peuvent être enrichis sans casser les anciens consommateurs si leurs schémas sont versionnés.

\subsection{État, événements et trace}
L'état d'exécution est $s=(D,K,M,B,C)$ : données locales $D$, contexte et connaissances accessibles $K$, mémoire durable $M$, budgets restants $B$ et capacités autorisées $C$. Une étape applique une transition
\[
\langle p,s\rangle \longrightarrow \langle p',s',r,e\rangle,
\]
où $p$ est le programme restant, $r$ le résultat partiel et $e$ l'événement de trace. La trace est une séquence append-only d'événements identifiés, horodatés, versionnés et chaînés par hash si l'exigence d'intégrité le requiert. Un hash signale un changement de contenu; à lui seul, il ne prouve ni la vérité de l'événement ni l'identité de son émetteur.

\section{Algèbre de programmes compositionnels}
\subsection{Syntaxe abstraite extensible}
\begin{definition}[Langage de composition]
Une grammaire de base est donnée par :
\begin{align*}
p ::= {} & \mathsf{Id} \mid \mathsf{Call}(c,e) \mid \mathsf{Seq}(p_1,\dots,p_n) \mid \mathsf{Par}(p_1,\dots,p_n) \\
&\mid \mathsf{If}(g,p_t,p_f) \mid \mathsf{Map}(p,e,n) \mid \mathsf{Fold}(p,e,a) \\
&\mid \mathsf{Loop}(g,p,I) \mid \mathsf{Choose}(p_1,\dots,p_n) \mid \mathsf{Verify}(p,v) \\
&\mid \mathsf{Observe}(q) \mid \mathsf{Probe}(Q) \mid \mathsf{Recover}(p,\rho) \\
&\mid \mathsf{Txn}(p,\kappa) \mid \mathsf{Checkpoint}(p,k) \mid \mathsf{Revoke}(d) \\
&\mid \mathsf{Compose}(m,\theta) \mid \mathsf{Transform}(p,\tau) \mid \mathsf{Fix}(x,p)\\
&\mid \mathsf{Select}(\pi,p_1,\dots,p_n) \mid \mathsf{Synthesize}(G,D) \\
&\mid \mathsf{Hibernate}(c) \mid \mathsf{Wake}(c,q) \mid \mathsf{Compress}(m,k) \\
&\mid \mathsf{Summarize}(m,k,\epsilon) \mid \mathsf{Restore}(k) \\
&\mid \mathsf{EmitWave}(w) \mid \mathsf{ReceiveWave}(w) \mid \mathsf{ResumeSearch}(k).
\end{align*}
Les formes non présentes dans un DSL existant sont des extensions théoriques à spécifier et implémenter, pas des capacités attribuées rétroactivement au démonstrateur. La grammaire doit être extensible par modules typés et versionnés, avec analyse syntaxique, validation sémantique, compilation et tests différentiels.
\end{definition}

\subsection{Composition séquentielle}
Pour $f:A\to B$ et $g:B\to C$, la composition $g\circ f:A\to C$ est bien typée si la sortie de $f$ correspond à l'entrée de $g$ et si les préconditions sont satisfaites. Au niveau des contrats, une règle suffisante est :
\[
\{P\}\ f\ \{Q\},\quad Q\Rightarrow P_g,\quad \{P_g\}\ g\ \{R\}
\quad\Longrightarrow\quad \{P\}\ g\circ f\ \{R\}.
\]
La condition $Q\Rightarrow P_g$ porte sur les états, les types, les unités, les effets et les hypothèses de contexte, pas seulement sur la forme de la donnée.

\subsection{Composition parallèle et ordre}
La composition parallèle ne peut être déclarée indépendante qu'après analyse des effets. Soient $R_i$ les zones lues, $W_i$ les zones écrites et $X_i$ les effets externes du composant $i$. Une condition conservatrice d'indépendance est
\[
W_i\cap(R_j\cup W_j)=\varnothing,\quad W_j\cap(R_i\cup W_i)=\varnothing,
\]
avec absence de conflit d'effets externes et absence d'interférence par ressources partagées. Sous ces conditions, et si les sémantiques locales sont déterministes sur l'état observé, les deux ordres peuvent être équivalents relativement à la projection d'observation choisie. Sans ces conditions, l'ordonnancement fait partie de la sémantique; permuter deux nœuds est une transformation non autorisée par défaut.

\subsection{Identité, associativité et lois de réécriture}
L'identité $\mathsf{Id}$ ne change ni l'état observable ni le résultat, hormis les événements de trace explicitement déclarés. La séquence est associative au niveau de la composition des fonctions partielles pures; pour les processus avec effets, l'équivalence dépend de la règle d'observation des traces, des frontières de transaction et de l'ordre des exceptions. Toute loi d'optimisation doit donc porter ses hypothèses et son domaine d'équivalence. L'optimiseur ne peut pas supprimer une étape d'audit, une vérification ou un effet simplement parce que sa valeur de sortie paraît inutilisée.

\subsection{Branchement, collection, réduction et récursion}
Un branchement exige un prédicat typé et une branche de repli explicitée. Un $\mathsf{Map}$ doit déclarer le type de chaque élément, la politique d'erreur par élément et la limite de ressources; un $\mathsf{Fold}$ doit déclarer le type de l'accumulateur et l'invariant de réduction. Une boucle doit avoir un invariant $I$, une condition d'arrêt et, pour prouver la terminaison, une fonction de rang $\mu$ à valeurs dans un ordre bien fondé, telle que chaque itération non terminale diminue strictement $\mu$. Les programmes récursifs sans preuve de terminaison peuvent être exécutés avec budget, mais leur terminaison ne doit pas être prétendue.

\section{Sémantique et preuves de correction compositionnelle}
\subsection{Sémantique dénotationnelle et opérationnelle}
Une composante pure peut être interprétée par une fonction partielle $\llbracket c\rrbracket:A\rightharpoonup B$. Un programme avec effets est plutôt une transition sur état et trace :
\[
\llbracket p\rrbracket: (s,x)\mapsto (s',y,\trace,w),
\]
avec résultat non garanti si la précondition est violée, une capacité manque ou une ressource s'épuise. Les fonctions partielles et transitions ne doivent pas être ramenées à la même abstraction lorsqu'il existe des écritures, appels distants, horloges, aléatoire, observation du monde ou intervention humaine.

\begin{proposition}[Correction de la composition séquentielle]
Supposons que les triples de Hoare $\{P\}f\{Q\}$ et $\{Q\}g\{R\}$ soient valides dans une sémantique commune, que l'exécution de $f$ termine sur toute entrée satisfaisant $P$, et qu'aucune interférence externe ne modifie l'état entre la postcondition $Q$ et l'évaluation de la précondition de $g$. Alors $\{P\}g\circ f\{R\}$ est valide.
\end{proposition}
\begin{proof}
Soit un état initial $s$ satisfaisant $P$. Par la validité du premier triple et la terminaison supposée, l'exécution de $f$ donne un état $s_1$ satisfaisant $Q$. Le second triple s'applique donc à $s_1$ et donne un état final satisfaisant $R$. Ainsi toute exécution du composé conforme aux hypothèses satisfait $R$.
\end{proof}
Cette preuve est mathématique et conditionnelle; elle n'est pas une preuve machine du compilateur AAM ni d'une implémentation particulière.

\begin{proposition}[Composition par contrat]
Soient $c_1,\ldots,c_n$ des composants dont les contrats sont prouvés dans un même système de types. Si chaque connexion transmet une sortie satisfaisant les préconditions du composant suivant, si toutes les capacités exigées sont disponibles et si les effets sont sérialisés conformément à leur contrat, alors la composition séquentielle satisfait la composition des postconditions et conserve la trace de chaque étape.
\end{proposition}
\begin{proof}
Par induction sur $n$. Le cas $n=1$ est le contrat du composant. Pour $n+1$, l'hypothèse d'induction fournit la postcondition de la composition des $n$ premiers. La condition de raccord rend cette postcondition suffisante pour la précondition du composant $n+1$; la proposition de correction séquentielle achève l'étape inductive.
\end{proof}

\begin{proposition}[Équivalence sûre d'une composition parallèle]
Supposons deux composants déterministes $f$ et $g$ dont les effets sont indépendants au sens de la section précédente, qui observent et modifient des parties disjointes de l'état, et dont les traces sont comparées après projection sur les événements pertinents au contrat. Alors $f$ puis $g$ et $g$ puis $f$ produisent le même état final projeté et la même valeur projetée.
\end{proposition}
\begin{proof}
Comme $f$ ne modifie aucune donnée lue ou écrite par $g$, l'application de $g$ ne dépend pas du fait que $f$ ait déjà exécuté; symétriquement pour $f$. Les mises à jour sur des composantes disjointes de l'état commutent. Les valeurs produites sont donc identiques et les états finaux diffèrent au plus sur l'ordre d'événements que la projection de trace a explicitement choisi d'ignorer.
\end{proof}

\subsection{Vérification indépendante et certificats}
Un certificat de composition contient au minimum : hash de la spécification source normalisée, versions de composants et du vérificateur, contrats utilisés, graphe de dépendances, configuration d'exécution, résultat, trace, hypothèses et verdict. Le compilateur et le vérificateur ne doivent pas partager silencieusement la même logique critique lorsque l'indépendance est une exigence. Une compilation réussie signifie que le programme a satisfait le contrôle de compilation annoncé; ce n'est pas une preuve de correction métier. Un hash lie des octets à une empreinte; il ne valide pas le contenu.


\section{Système de types, règles de jugement et sémantique exécutable}
\subsection{Types et environnements}
Les types du noyau sont générés par la grammaire
\begin{align*}
\tau ::= {}& \mathbf{Unit}\mid\mathbf{Bool}\mid\mathbf{Int}\mid\mathbf{Real}\mid\mathbf{String}\mid\mathbf{Any}\\
&\mid \mathbf{List}(\tau)\mid \mathbf{Map}(\tau,\tau)\mid \mathbf{Pair}(\tau,\tau)\\
&\mid \mathbf{Sum}(\tau,\tau)\mid \mathbf{Ref}(\tau)\mid \mathbf{Record}(S).
\end{align*}
Les enregistrements $S$ ont un schéma versionné. Les types numériques peuvent être décorés d'une dimension d'unité; les références peuvent être décorées par leur portée et leur version. Le type $\mathbf{Any}$ n'est permis qu'à la frontière d'entrée/sortie opaque et déclenche un contrôle dynamique avant toute opération qui demande un type concret.

$\Gamma$ est l'environnement des variables, $\Sigma$ le registre versionné des signatures de composants, et $K$ l'ensemble des capacités d'exécution. Un jugement d'expression est $\Gamma\vdash e:\tau$. Un jugement de programme est
\[
\Gamma;\Sigma;K\vdash p: A\Rightarrow B\;!\;E\;[P,Q]\;\text{cost}\leq r,
\]
ce qui signifie : sous les variables disponibles, les signatures versionnées et les capacités $K$, le programme transforme une entrée $A$ en résultat $B$, peut produire les effets $E$, est soumis à la précondition $P$, vise la postcondition $Q$ et déclare une borne de ressources $r$. Le jugement n'affirme ces propriétés que si les contrats des composants référencés sont eux-mêmes valides à leur niveau de preuve déclaré.

\subsection{Règles de typage des expressions}
Les littéraux sont typés par leur type concret. Si $x:\tau\in\Gamma$, alors $\Gamma\vdash x:\tau$. Pour les listes, si chaque élément $e_i$ a le type $\tau$, la liste a le type $\mathbf{List}(\tau)$. Pour un opérateur pur $op$ avec signature $\tau_1\times\cdots\times\tau_n\to\tau$, les prémisses $\Gamma\vdash e_i:\tau_i$ donnent $\Gamma\vdash op(e_1,\ldots,e_n):\tau$. La règle est refusée si un opérateur n'est pas dans le registre ou si les types ne satisfont pas la signature. Une coercition éventuelle est un opérateur explicite avec sa propre sémantique, jamais une conversion silencieuse.

\subsection{Règles de typage des programmes}
\paragraph{Identité.}
\[
\frac{}{\Gamma;\Sigma;K\vdash \mathsf{Id}:A\Rightarrow A\;!\;\varnothing\;[\mathsf{true},\mathsf{true}]}.
\]
\paragraph{Appel.} Si $\Sigma(c)=(A,B,P_c,Q_c,E_c,r_c)$, $\Gamma\vdash e:A$, $K$ autorise $E_c$, et $P\Rightarrow P_c(e)$, alors
\[
\frac{\Sigma(c)=(A,B,P_c,Q_c,E_c,r_c)\quad\Gamma\vdash e:A\quad E_c\subseteq K\quad P\Rightarrow P_c(e)}
{\Gamma;\Sigma;K\vdash \mathsf{Call}(c,e):A\Rightarrow B\;!\;E_c\;[P,Q_c]}.
\]
La postcondition de l'appel est instanciée par les entrées et sorties du composant. Une capacité manquante rend l'appel non admissible pour cette exécution; elle ne transforme pas l'effet en pur.
\paragraph{Séquence.} Si $p_1:A\Rightarrow B\;!E_1[P,Q_1]$ et $p_2:B\Rightarrow C\;!E_2[P_2,Q_2]$ sont typés et $Q_1\Rightarrow P_2$, alors
\[
\frac{\Gamma;\Sigma;K\vdash p_1:A\Rightarrow B\;!E_1[P,Q_1]\quad
\Gamma;\Sigma;K\vdash p_2:B\Rightarrow C\;!E_2[P_2,Q_2]\quad Q_1\Rightarrow P_2}
{\Gamma;\Sigma;K\vdash \mathsf{Seq}(p_1,p_2):A\Rightarrow C\;!(E_1\cup E_2)[P,Q_2]}.
\]
\paragraph{Branchement.} Si $g$ est booléen et les deux branches ont les mêmes types d'entrée/sortie et une postcondition commune $Q$, le branchement a l'union de leurs effets. La postcondition s'applique à la branche sélectionnée, pas à la branche non exécutée.
\paragraph{Map.} Si $p:A\Rightarrow B$ et $e:\mathbf{List}(A)$, alors $\mathsf{Map}(p,e,n)$ produit $\mathbf{List}(B)$ uniquement si la politique sur les erreurs, le maximum $n$ et la consommation de ressources sont déclarés. En politique $\mathsf{StopOnError}$, le résultat global n'est pas $\mathsf{Success}$ si une itération échoue; en politique $\mathsf{Collect}$, le type de sortie doit être un type somme qui préserve séparément résultats et erreurs.
\paragraph{Vérification.} $\mathsf{Verify}(p,v)$ ne transforme pas arbitrairement un résultat : elle produit un type de résultat distinguant validation, rejet et vérification indéterminée. La promotion épistémique dépend ensuite du contrat du vérificateur.
\paragraph{Récupération.} $\mathsf{Recover}(p,q)$ est typé si la sortie de $q$ respecte le contrat minimal exigé par l'appelant lorsque $p$ échoue. Une récupération ne peut être considérée comme sûre que si les effets de $p$ sont purs, transactionnels, idempotents ou compensables, et que sa politique de restauration est spécifiée.
\paragraph{Boucle.} $\mathsf{Loop}(g,p,I)$ exige un garde booléen et un invariant $I$. Pour prétendre à la terminaison, le contrat ajoute une fonction de rang $\mu$ qui décroît strictement dans un ordre bien fondé à chaque itération non terminale.
\paragraph{Parallélisme.} $\mathsf{Par}(p_1,p_2)$ n'est admissible que si les effets se composent, que l'interférence est exclue ou résolue par une politique formelle, et que le mécanisme de fusion des sorties est défini. L'absence de preuve d'indépendance conduit à conserver un ordre séquentiel ou à demander un raffinement des contrats.

\subsection{Sémantique à grands pas}
La relation d'exécution est
\[
\langle p,s\rangle\Downarrow\langle y,s',\trace,c\rangle,
\]
où $y$ est un résultat typé (succès, partiel, non résolu, rejet, échec ou annulation), $s'$ l'état final, $\trace$ la suite des événements et $c$ le coût cumulé. La trace inclut les appels de composants, branches, probes, validations, effets, erreurs, récupérations et décisions de budget.

\paragraph{Appel.} L'évaluation d'un appel suit les phases strictes : résolution du composant et de sa version; résolution des expressions; validation du type d'entrée et des préconditions; vérification des capacités; exécution; validation du type de sortie et de la postcondition; écriture de l'événement de résultat. Si une phase échoue, la sémantique renvoie un statut explicite et conserve les événements déjà survenus. Les effets externes déjà réalisés ne sont pas réputés annulés.

\paragraph{Séquence.} Si $p_1$ produit $\mathsf{Success}(v)$ et un état $s_1$ satisfaisant le contrat de raccord, alors $p_2$ s'exécute sur $s_1$. Si $p_1$ renvoie un statut d'échec ou de rejet, la politique par défaut propage ce statut et n'exécute pas $p_2$. Une politique de continuation est possible, mais doit être indiquée dans l'AST et spécifier comment les résultats partiels sont agrégés.

\paragraph{Branchement.} Le garde est évalué à partir de l'état courant, son résultat est consigné dans la trace, puis une seule branche est exécutée. Une évaluation de garde indéterminée ne choisit pas une branche par défaut : elle produit un statut non résolu ou appelle une procédure de clarification explicitement déclarée.

\paragraph{Map et Fold.} L'ordre, les limites, la politique d'erreur et l'agrégation sont partie intégrante de la sémantique. L'environnement local d'un élément est restauré après son traitement, sauf si le contrat déclare expressément une réduction d'état. Un fold doit préciser son accumulateur initial et sa loi de réduction; l'associativité de la réduction n'est pas présumée.

\paragraph{Vérification.} Une vérification produit l'un des statuts $\mathsf{VALID}$, $\mathsf{INVALID}$ ou $\mathsf{UNKNOWN}$, chacun accompagné d'un certificat, d'un contre-exemple ou d'une raison d'insuffisance. Seul le statut $\mathsf{VALID}$, dans la portée déclarée et après le contrôle d'intégrité du certificat, peut servir de preuve positive pour le contrat concerné.

\paragraph{Récupération.} Le runtime sauvegarde le sous-état contrôlé avant d'exécuter le corps. Si le corps échoue, les mutations réversibles sont annulées ou compensées, l'événement de récupération est tracé, puis le repli est exécuté. Si un effet irréversible est en cours de réconciliation, le statut global doit le signaler; la restauration mémoire seule ne suffit pas.

\subsection{Progrès, préservation et sûreté de typage}
\begin{theorem}[Sûreté du sous-langage pur]
Considérons le fragment fini composé de $\mathsf{Id}$, $\mathsf{Call}$, $\mathsf{Seq}$, $\mathsf{If}$, $\mathsf{Map}$ sur une liste finie et $\mathsf{Verify}$, à l'exclusion de $\mathsf{Loop}$, $\mathsf{Fix}$, du parallélisme et de tout effet externe. Supposons un registre fini de composants purs, totaux et conformes à leurs signatures, des opérateurs purs totaux sur leurs domaines typés, des gardes et vérificateurs totaux, ainsi qu'une grammaire correctement vérifiée. Si $\Gamma\vdash p:A\Rightarrow B$ et si l'entrée $x$ a le type $A$, alors l'exécution termine avec un résultat de type $B$ ou avec une erreur typée explicitement représentée par la sémantique; elle ne se bloque pas sur une opération de type inconnu.
\end{theorem}
\begin{proof}
Par induction structurale sur $p$. L'identité renvoie son entrée. L'appel vérifie son entrée, puis le composant total fournit une sortie du type annoncé. Pour la séquence, le type de sortie du premier programme fournit le type requis au second. Le branchement choisit une branche dont le résultat est du type commun. Le map traite un nombre fini d'éléments et applique l'hypothèse d'induction à chacun. Le vérificateur est total et renvoie l'un des cas du type-somme déclaré. Les opérations de type invalide sont arrêtées par la vérification et deviennent des erreurs typées; aucune règle du fragment ne crée une opération non représentée par la sémantique.
\end{proof}
Le théorème ne couvre pas les effets externes, les programmes infinis, ni les composants dont l'implémentation viole leur contrat. Ces cas exigent les règles d'effets, de terminaison et de résilience qui suivent.

\begin{theorem}[Terminaison conditionnelle d'une boucle]
Supposons que l'invariant $I$ soit vrai avant la boucle, qu'il soit préservé par chaque exécution non terminale du corps, que le garde et le corps terminent à chaque itération, et qu'une mesure $\mu$ décroisse strictement à chaque itération tout en appartenant à un ordre bien fondé. Alors la boucle termine et laisse un état satisfaisant l'invariant et la négation du garde.
\end{theorem}
\begin{proof}
Un ordre bien fondé ne contient aucune suite infinie strictement décroissante; la mesure ne peut donc décroître indéfiniment. À l'itération finale, le garde est faux. L'invariant est vrai par induction sur le nombre fini d'itérations et demeure vrai à la sortie.
\end{proof}

\subsection{Lois de transformation et équivalence comportementale}
Soient $p$ et $q$ deux programmes. Leur équivalence dépend d'une observation contractuelle $Obs$ :
\[
p\simeq_{Obs}q\iff\forall x\in D,\ Obs(\llbracket p\rrbracket(x))=Obs(\llbracket q\rrbracket(x)).
\]
Le domaine $D$, les versions de composants, l'environnement et les effets observés sont partie intégrante de cette affirmation. Une équivalence limitée à des entrées finies testées est notée $\mathsf{EquivalentOn}(D_{test})$, pas $\simeq_{Obs}$ universelle. Les transformations automatiques exigent une preuve d'équivalence dans le sous-langage pris en charge ou un rapport de tests bornés accompagné de sa portée.

\begin{proposition}[Préservation d'une transformation certifiée]
Si une transformation $T$ fournit un certificat valide démontrant $p\simeq_{Obs}T(p)$ sous des hypothèses $H$, et si le contexte d'exécution satisfait encore $H$, remplacer $p$ par $T(p)$ préserve les propriétés observables définies par $Obs$ dans ce contexte.
\end{proposition}
\begin{proof}
Par définition de l'équivalence observée, les résultats observables de $p$ et $T(p)$ sont égaux pour toutes les entrées du domaine $D$ sous $H$. Le remplacement conserve donc les propriétés exprimées uniquement par $Obs$ sur ce domaine. Aucun comportement extérieur à $D$ ni propriété absente de $Obs$ n'est déduit.
\end{proof}

\subsection{Correspondance interpréteur/compilateur}
Un compilateur est correct si l'exécution du programme compilé et celle de son AST source sont reliées par une relation de simulation préservant les résultats, les effets autorisés et la trace contractuelle. Le but formel est :
\[
\begin{aligned}
\forall p,s,\quad \mathsf{Valid}(p)\Rightarrow {}&
\operatorname{Obs}(\llbracket p\rrbracket_{\mathrm{interp}}(s))\\
={}&\operatorname{Obs}(\llbracket \operatorname{Compile}(p)\rrbracket_{\mathrm{run}}(s)).
\end{aligned}
\]
Cette formule est une obligation de preuve, pas un résultat déclaré acquis dans le présent dossier. Elle se ferme par (i) sémantique formelle des nœuds; (ii) preuve par induction sur l'AST pour l'émission du code; (iii) tests différentiels générés; (iv) tests de rejet négatifs; (v) contrôle d'empreintes et de versions. Les tests différentiels ne remplacent pas la preuve générale mais peuvent détecter un contre-exemple à une version du compilateur.

\section{Transactions, effets externes et sémantique de récupération}
\subsection{État transactionnel}
Une transaction est $\mathsf{Txn}(p,\kappa)$ où $\kappa$ définit les frontières de commit, rollback et compensation. Son état comprend un journal de mutations préparées, les effets confirmés et les clés d'idempotence. Le commit est autorisé lorsque les préconditions/postconditions du programme et les contrôles requis réussissent. En cas d'échec avant commit, les mutations locales transactionnelles sont annulées. Les effets externes non transactionnels doivent être inscrits au journal comme irréversibles ou soumis à une action compensatoire.

\subsection{Reprise et sémantique de livraison}
Après interruption, un checkpoint est accepté seulement si ses hashes, versions, schémas et séquence d'événements sont cohérents. Une opération externe avec identifiant stable et accusé de réception connu peut être marquée confirmée; une opération non idempotente dont l'issue est inconnue doit être réconciliée avant nouvel envoi. Les garanties « au moins une fois » et « au plus une fois » ne sont pas équivalentes; « exactement une fois » requiert une déduplication durable ou un protocole coopérant avec le système externe. Il est interdit de revendiquer exactement une fois sur la seule base d'un checkpoint local.

\begin{proposition}[Reprise idempotente]
Supposons qu'une étape $f$ ait un identifiant d'opération persistant, que le destinataire conserve durablement les clés déjà exécutées, et qu'un replay avec la même clé renvoie le résultat mémorisé sans répéter l'effet. Alors la reprise après crash peut réexécuter $f$ avec la même clé sans dupliquer l'effet externe, sous réserve que la conservation de la clé et de son résultat soit atomique côté destinataire.
\end{proposition}
\begin{proof}
Lors du premier traitement, le destinataire enregistre le couple clé/résultat avec l'effet. Lors du replay, la même clé est reconnue avant l'exécution de l'effet; le résultat mémorisé est renvoyé. La condition d'atomicité empêche un état dans lequel l'effet est réalisé mais la clé est perdue. Ainsi le replay n'ajoute pas de seconde occurrence de l'effet dans les hypothèses données.
\end{proof}

\subsection{Cohérence des checkpoints}
Le checkpoint doit contenir ou référencer : id et version de la composition; état sérialisé; état des variables; événements de trace et dernier hash confirmé; opérations externes avec leur état de réconciliation; versions des composants; budget; artefacts et dépendances. Le redémarrage rejoue uniquement les étapes dont les effets le permettent et dont les préconditions sont encore vraies. Les événements non confirmés sont recontrôlés; une donnée périmée déclenche l'invalidation des descendants.

\section{Intégrité des traces et portée des empreintes}
La trace est chaînée par $h_0=H(\text{seed})$ et $h_i=H(\operatorname{Canon}(e_i,h_{i-1}))$, où $H$ est un hash cryptographique et $\operatorname{Canon}$ une sérialisation canonique versionnée. Toute modification d'un événement ou de son prédécesseur entraîne une incohérence des hashes suivants, à moins d'une collision ou d'une réécriture complète du suffixe. La chaîne n'est pas une signature : pour attribuer la trace à un acteur, il faut une clé, une signature ou une attestation indépendante. Le manifeste relie les artefacts, versions et paramètres mais ne certifie pas à lui seul la correction sémantique.

\section{Limites logiques générales et obligations d'extension}
Un ensemble fini d'exemples peut être expliqué par plusieurs programmes distincts. L'identifiabilité dépend du domaine, des contraintes, des interventions et de la classe de candidats; la composition ne supprime pas cette propriété logique, elle rend explicites les briques et les hypothèses utilisées. Quand les données ne distinguent pas les candidats, le système doit conserver le sous-ensemble compatible et rechercher une observation qui les sépare. Si la grammaire de candidats ne représente pas la fonction recherchée, elle doit être enrichie par une nouvelle forme de programme, avec preuve de typage, analyse d'effets, tests et trace.

Pour qu'une extension s'intègre, elle doit clôturer un dossier contenant : définition mathématique; syntaxe/AST; règles de typage; sémantique opérationnelle; effets/capacités; préconditions/postconditions; argument de correction; conditions de terminaison; coûts et limites d'exécution; stratégie de panne/reprise; tests positifs et négatifs; version et compatibilité; certificats/provenance; comparaison avec les voies existantes. Une entrée de backlog seule ne satisfait pas cette obligation; elle doit être transformée en sous-projet jusqu'au niveau de formalisation et d'implémentation requis pour son contrat d'intégration.


\section{Effets, capacités, droits et sûreté de l'exécution}
\subsection{Système d'effets}
Chaque composant annonce ses effets dans un ensemble typé. Un exemple de vocabulaire d'effets est
\[
\begin{aligned}
E=\{&\mathsf{Read}(p),\mathsf{Write}(p),\mathsf{Network}(h),\mathsf{Clock},\\
&\mathsf{Random},\mathsf{Observe}(t),\mathsf{Execute}(k)\}.
\end{aligned}
\]
Les effets sensibles exigent une capacité détenue, un périmètre, une politique et une trace. L'absence de capacité est une condition de refus ou de recherche d'une voie alternative autorisée, et non une permission implicite.

\subsection{Frontière entre code et données}
La source compositionnelle est une représentation de programme, pas du code arbitraire déguisé en donnée. Le compilateur analyse le schéma, résout les identifiants dans un registre, contrôle les expressions, puis compile le programme validé. Les littéraux, noms de composants, attributs et chemins ne sont pas équivalents à du code exécutable libre. Toute extension de grammaire qui introduit un nouveau moyen d'exécution doit déclarer son modèle de menace, ses capacités requises et ses tests adversariaux.

\subsection{Confinement sans fermeture évolutive}
La sécurité d'exécution n'exige pas que le langage demeure figé. Elle exige que chaque nouvelle construction passe par le cycle : spécification, typage, analyse des effets, tests négatifs, revue du compilateur, vérification du manifeste, contrôle de compatibilité, déploiement versionné et possibilité de retour arrière. Si une tâche dépasse la grammaire actuelle, le travail correct consiste à spécifier puis ajouter l'abstraction manquante avec ses obligations de sûreté, plutôt qu'à laisser le besoin hors du modèle.

\section{Résilience : erreurs, reprise, transactions et continuité}
\subsection{Taxonomie des défaillances}
Les erreurs sont classées au minimum en : entrée invalide, précondition non satisfaite, contrat de sortie rompu, composant absent, permission insuffisante, indisponibilité temporaire, timeout, manque de budget, état conflictuel, corruption d'artefact, échec de vérification, résultat indéterminé, et erreur de l'orchestrateur. Les catégories déterminent une action de reprise; une exception générique ne suffit pas à sélectionner une stratégie.

\subsection{Politiques de reprise}
Un composant peut déclarer une politique parmi : aucun nouvel essai, nouvel essai borné, nouvel essai avec backoff, stratégie alternative, récupération depuis checkpoint, compensation de l'effet, retour arrière de transaction, continuation partielle ou arrêt contrôlé. La reprise n'est sûre que si elle est idempotente, dédupliquée par identifiant de commande, ou compensée. Il est incorrect de traiter automatiquement toutes les erreurs comme transitoires.

\subsection{Transactions et effets externes}
Une transaction locale peut offrir atomicité si le backend la garantit. Pour une suite de systèmes externes qui ne partagent pas une transaction distribuée, le modèle doit utiliser une orchestration de compensation (saga) ou déclarer explicitement les effets irréversibles. Un simple rollback de mémoire ne peut pas annuler un courriel envoyé, une commande physique ou une écriture distante déjà confirmée.

\subsection{Reprise après interruption}
Un checkpoint sauvegarde état sérialisé, version de schéma, registre de composants, derniers événements confirmés, budgets, empreintes d'artefacts et identifiants d'opérations externes. Au redémarrage, l'exécuteur valide le checkpoint, reconstruit les dépendances et détermine quels effets ont été confirmés. Les actions non décidables ne sont pas rejouées aveuglément : elles entrent dans un état de réconciliation. La continuité de travail exige une opération atomique de sauvegarde, des copies versionnées et un journal append-only.

\subsection{Disponibilité dégradée et résultat partiel}
Quand un composant devient indisponible, l'orchestrateur recherche une voie alternative compatible avec le contrat et les capacités disponibles, réutilise des résultats certifiés encore valides ou renvoie un résultat partiel explicitement étiqueté. La résilience préserve autant que possible la fonction; elle ne maquille pas la perte de garanties. Un fallback n'est accepté que si son contrat couvre la demande ou si la sortie expose précisément l'écart.

\begin{invariant}[Échec explicite et non-promotion]
Aucun résultat marqué partiel, inconnu, périmé, non vérifié ou interrompu ne peut être converti en résultat certifié sans une nouvelle étape de vérification dont les préconditions sont satisfaites.
\end{invariant}

\section{Composition de données, mémoire et provenance}
Chaque valeur produite est rattachée à : composant, version, paramètres normalisés, hash des entrées pertinentes, contexte, temps, effets, artefacts de sortie, statut et dépendances. Un graphe de provenance $G=(V,E)$ lie une sortie à ses prémisses et à ses étapes de transformation. Les arêtes distinguent au moins dépendance de valeur, dépendance de contrôle, dépendance de preuve et dépendance d'environnement.

\begin{proposition}[Propagation conservatrice d'invalidation]
Soit un graphe de dépendances orienté $G$ où $u\to v$ signifie que la validité de $v$ dépend de $u$. Si un nœud $u$ est révoqué, l'ensemble des descendants accessibles de $u$ est marqué à revalider ou invalider avant qu'un de ces nœuds soit réutilisé comme prémisse certifiée.
\end{proposition}
\begin{proof}
Par définition de l'arête, la garantie de validité de chaque descendant dépend au moins indirectement de $u$. Après révocation de $u$, la conjonction d'hypothèses qui justifiait cette garantie n'est plus établie. En parcourant le graphe dans l'ordre topologique, chaque descendant est donc placé en attente de revalidation; une promotion ultérieure nécessite une preuve alternative ou la reconstruction de ses dépendances valides.
\end{proof}
Cette règle suppose que les dépendances sont exhaustives et correctement enregistrées. En présence de cycles, les composants fortement connexes sont traités comme des unités de révision; toute dépendance cachée non enregistrée constitue une violation de l'invariant de provenance.

\section{Synthèse, recherche et transformation de programmes}
\subsection{La composition comme objet de recherche}
Le système ne doit pas seulement exécuter des compositions rédigées à l'avance. Il doit pouvoir rechercher une composition $p$ dans une grammaire $\mathcal{G}$ telle que, pour un jeu d'exemples d'entraînement $D_{tr}$, $p$ satisfasse les contrats et minimise un objectif vectoriel :
\[
J(p)=\big(L_{ex}(p),L_{contraintes}(p),L_{transfert}(p),C(p),R_{effets}(p),U(p)\big),
\]
où les termes représentent respectivement l'écart aux exemples, les violations de contraintes, la perte de transfert, le coût, le risque lié aux effets et l'incertitude estimée. Les termes incompatibles restent séparés jusqu'à ce que la politique de décision justifie un arbitrage. Une pénalité basse ou une solution courte n'établit pas sa correction.

\subsection{Boucle de synthèse}
La boucle de synthèse comporte : formalisation du besoin; définition des observables et de l'oracle de test; génération de programmes dans une grammaire typée; normalisation des formes équivalentes; vérification statique; exécution sur cas entraînement; recherche de contre-exemples; raffinement; tests holdout; tests perturbés; analyse d'effets; certification conditionnelle; archivage du programme et de sa provenance. Le jeu holdout reste séparé du générateur et du normaliseur; les réponses attendues ne peuvent pas devenir des variables d'entrée de synthèse.

\subsection{Décomposition et réutilisation hiérarchique}
Un programme complexe peut être factorisé en sous-compositions nommées avec contrat, version et coût, puis recomposé à un niveau supérieur. Une décomposition utile doit conserver le comportement observable annoncé. La factorisation n'est pas correcte simplement parce qu'elle produit du code plus court; elle doit préserver les postconditions, les effets, l'ordre des événements et les garanties de ressources exigées.

\subsection{Composition évolutive et auto-transformation}
Un programme peut proposer des transformations de lui-même ou de ses sous-compositions : remplacement d'un composant, factorisation, fusion, parallélisation, spécialisation, généralisation, changement de stratégie de recherche, ajout d'une étape de validation, ou réorganisation des budgets. Ces transformations sont des candidats. Pour être mises en service elles doivent satisfaire un contrat de transformation, produire un diff, passer les vérifications et tests, être comparées sur un ensemble de régression, et être livrées avec la possibilité de rollback. La version actuellement active reste identifiable, et une dégradation détectée entraîne la restauration du dernier état approuvé.

\subsection{Généralisation compositionnelle}
Le test de compositionnalité doit séparer les briques vues et les compositions vues. Un split pertinent tient les primitives disponibles mais réserve des combinaisons inédites; des splits plus stricts réservent aussi des familles de primitives et des contextes. Les métriques comprennent l'exactitude sur composition inédite, la profondeur atteinte, la sensibilité à l'ordre, la robustesse aux noms permutés, la réutilisation correcte de contrats, le taux de violation de contraintes et le coût de recherche. Un test qui contient des gabarits presque identiques au jeu d'entraînement mesure surtout la mémorisation.

\section{Programmation compositionnelle multimodale et sémantique}
Le noyau traite uniformément les contrats, mais garde les types de domaine distincts. Une composition de langage peut lier tokenisation déterministe, parseur, analyse de référents, interprétation de portée, mémoire contextuelle, résolution d'ambiguïté, contrôle d'entailment et vérificateur. Une composition numérique peut lier parseur d'expression, système d'unités, simplificateur, solveur et vérificateur de preuve. Une composition d'observation peut lier collecte autorisée, empreintes multi-vues, modèle de transition, sondes et comparaison de prédictions. Tous sont des plans de composition différents sur un même cadre de contrats; aucune couche linguistique ne bénéficie d'un passe-droit épistémique.

Les idées multirésolutions sont transférables sous forme de vues ou projections hiérarchiques sur les données (octets, tokens, syntagmes, structure, sémantique, trajectoire) uniquement si chaque vue apporte une utilité testable et conserve sa provenance. Aucune architecture de FulleGenBrain ne dépend d'un modèle génératif de langage ou de poids d'un tel modèle. Les capacités de formulation et de découverte proviennent de représentations explicites, de données observées, de grammaires de programmes, de recherche, d'analogies calculables, de moteurs spécialisés et de composition.

\section{Programmation par composition comme pilier de FulleGenBrain\_2}
La programmation compositionnelle n'est pas un service externe dans FulleGenBrain\_2 : elle est le mécanisme structurel par lequel le système assemble ses voies cognitives, transforme ses plans de calcul, apprend de nouveaux comportements et répare ses propres chaînes d'action. Les composants fondamentaux sont des compétences (observer, parser, calculer, chercher, mesurer, induire, simuler, vérifier), des méthodes (stratégies réutilisables), des règles/invariants (contraintes transversales) et des adaptateurs de domaine. Les compositions sont des objets typés enregistrés en mémoire et exécutables par un runtime gouverné.

\subsection{Architecture à trois plans}
\begin{enumerate}
\item \textbf{Plan de définition :} catalogue versionné, contrats, DSL/AST, types, effets, budget, politiques d'erreur, schémas de provenance et tests.
\item \textbf{Plan d'exécution :} interpréteur de référence, compilateur, ordonnanceur, contrôleur d'état, sandbox par capacités, journal et checkpoints.
\item \textbf{Plan d'évolution :} induction de candidats, découverte de sous-programmes, synthèse, test différentiel, évaluation sur holdout, approbation ou retour arrière.
\end{enumerate}
Les plans sont séparés mais reliés par un identifiant de composition et un contrat commun. L'évolution peut proposer les changements; seul le chemin de validation déterminé peut activer une version nouvelle.

\subsection{Répartition des responsabilités}
Le contrôleur décide quel objectif poursuivre et quand s'arrêter; il n'exécute pas implicitement chaque compétence. Le synthétiseur propose des arbres de composition; il ne les certifie pas. Le compilateur contrôle que le programme est bien formé et produit un artefact reproductible; il ne certifie pas le résultat métier. Le runtime exécute uniquement les composants résolus par registre et les effets autorisés. Le vérificateur évalue les critères indépendamment de l'heuristique de recherche. La mémoire conserve artefacts, essais, dépendances et versions; elle ne convertit pas une sortie en vérité. Le gestionnaire de résilience observe l'état de santé des composants et choisit les reprises déclarées.

\subsection{Interfaces minimales des blocs}
Tout bloc doit déclarer : identifiant/version; entrée/sortie typées; préconditions/postconditions; effets et capacités; coût prévu et limite; exceptions/résultats partiels possibles; stratégie de compensation ou raison de non-compensabilité; provenance; tests unitaires; vérificateur; règles de compatibilité; comportement de révocation. Une ancienne version demeure disponible tant qu'elle est référencée par une composition active, un checkpoint ou une preuve historique.

\section{Protocole de validation exhaustif et évolutif}
La validation doit être une matrice de contrats, de composition et de tâches réelles, pas une série unique de démonstrations.
\begin{longtable}{p{0.20\textwidth}p{0.36\textwidth}p{0.34\textwidth}}
\toprule Niveau & Test exigé & Preuve conservée \\
\midrule
Composant & typage, pré/postconditions, valeurs frontières, erreurs, effets & logs, tests, versions, coverage et hashes \\
Connexion & compatibilité de types, unités, contextes, provenance & rapport de connexions refusées/acceptées \\
Composition & séquence, branche, map/fold, parallèle, boucle, reprise, transaction & trace et état final reproductible \\
Transformation & sémantique avant/après, diffs, régression, fallback & rapport de comparaison et décision d'activation \\
Résilience & panne injectée, interruption, timeout, restauration, doublons & journal de récupération et effets réconciliés \\
Synthèse & exemples, contre-exemples, ambiguïté, holdout, OOD & candidats, budget, oracle séparé, métriques \\
Sécurité & privilège minimal, identifiant non enregistré, chemin invalide, exécution indirecte & refus explicites et test red-team \\
Intégration & charge, contention, mémoire, dérivations révoquées, versions mixtes & profils, incidents, reprise et résultat final \\
\bottomrule
\end{longtable}

Les portes de sortie sont indépendantes : syntaxe/compilation, typage, contrat, sécurité des effets, intégrité de trace, qualité métier, généralisation et coût. Le résultat global rapporte le vecteur de scores et les échecs par catégorie. Les cas non exécutés sont marqués \code{NOT_RUN}; un budget épuisé signifie que l'essai est incomplet, pas que la capacité est absente. Tout cas sous-déterminé conserve plusieurs modèles ou demande une observation discriminante.

\section{Obligations théoriques restantes et stratégie de fermeture}
Une théorie complète et harmonieuse n'est pas acquise par la seule présence des définitions. Les obligations suivantes doivent être refermées progressivement par une formalisation et des artefacts testables :
\begin{enumerate}
\item Définir une logique de contrats et une théorie des types suffisamment expressive pour les unités, les contextes et les effets.
\item Spécifier la sémantique complète de chaque nœud DSL et démontrer la correspondance entre interpréteur et compilateur.
\item Formaliser concurrence, annulation, équité de l'ordonnancement et partage des ressources.
\item Formaliser transactions locales, sagas, opérations idempotentes et réconciliation post-crash.
\item Définir les règles de compatibilité/versionnement des composants et des compositions persistées.
\item Construire des preuves machine pour les lemmes prioritaires et un noyau de vérification minimal.
\item Mesurer empiriquement la capacité de synthèse et de généralisation sur des splits compositionalement stricts.
\item Fermer la chaîne d'invalidation pour les dépendances causales, sémantiques et d'environnement qui ne sont pas directement des dépendances de données.
\end{enumerate}
Ces obligations sont un programme de travail positif, pas une prédiction d'échec. Si un type, un opérateur ou un cas d'usage nécessaire n'est pas pris en charge, le processus prescrit est de le spécifier, le formaliser, l'implémenter, le vérifier et l'intégrer, en maintenant un statut épistémique précis durant le travail.

\section{Synthèse des garanties et des statuts}
\begin{longtable}{p{0.30\textwidth}p{0.26\textwidth}p{0.35\textwidth}}
\toprule Propriété & Statut dans ce document & Portée exacte \\
\midrule
Correction séquentielle sous contrats & PROVED-UNDER-ASSUMPTIONS & preuve mathématique fournie, hypothèses explicites \\
Indépendance parallèle & PROVED-UNDER-ASSUMPTIONS & uniquement pour effets et lectures/écritures disjoints \\
Propagation d'invalidation & PROVED-UNDER-ASSUMPTIONS & graphe exhaustif, descendants réévalués \\
DSL AAM avec 4 nœuds de base & SOURCE-OBSERVED & artefact de démonstration analysé \\
Compilateur AAM interdit les fragments source libres dans le DSL étudié & SOURCE-OBSERVED & implémentation présente, revue non exhaustive de sécurité \\
Équivalence compilateur/interpréteur complète & EMPIRICAL-TO-TEST & à couvrir par tests différentiels et preuve de traduction \\
Synthèse généralisant à des compositions inédites & EMPIRICAL-TO-TEST & aucune généralisation globale déduite du formalisme \\
FulleGenBrain\_2 intégré de bout en bout & À CONSTRUIRE/ÉVALUER & ce document définit l'architecture cible, pas une exécution intégrée \\
\bottomrule
\end{longtable}

\section{Conclusion}
La programmation compositionnelle devient un pilier lorsqu'elle fournit le langage de structure des capacités, le système de contrats, la sémantique d'exécution, la voie de synthèse et le mécanisme de transformation contrôlée. Sa solidité dépend de lois explicites, de préconditions compatibles, d'effets analysés, de preuves traçables, de tests adversariaux et de la capacité à reprendre après panne sans masquer les écarts. La meilleure architecture n'est pas celle qui interdit d'évoluer : elle est celle où chaque évolution est représentée, construite, examinée, testée, traçée et, si nécessaire, annulée proprement.

\end{document}
